\section*{Conclusão}

\begin{frame}{Fechando o ciclo}
  \begin{itemize}
    \item O semicondutor é um cristal comum cuja condutividade nós aprendemos
          a \termo{controlar}, e quem define esse controle é a banda proibida.
    \item A \termo{pureza de 11N} é o que permite que a dopagem intencional,
          na casa de ppm, signifique alguma coisa dentro do cristal.
    \item Uma sequência de reações bem conhecidas leva do quartzo ao wafer,
          passando por redução, destilação do \qui{SiHCl_3}, cristalização e
          deposição.
    \item O preço acompanha a \termo{pureza e o controle de processo}, e não a
          quantidade de material, e por isso o valor sobe mil vezes ao longo
          do caminho.
  \end{itemize}
\end{frame}

\begin{frame}[plain]
  \vfill
  \centering
  {\Large Obrigado! \\[0.3cm] Perguntas?}
  \vfill
\end{frame}
